\documentclass[11pt]{article}
\usepackage[margin=1in]{geometry}
\usepackage{hyperref}
\usepackage{graphicx}

% ── Paragraph spacing ────────────────────────────────────────────────────────
\usepackage{parskip}

% ── Page-break quality ───────────────────────────────────────────────────────
\widowpenalty=10000
\clubpenalty=10000
\raggedbottom

% ── Suppress automatic section numbering ─────────────────────────────────────
\setcounter{secnumdepth}{-1}

% ── Pandoc-generated table support ───────────────────────────────────────────
\usepackage{longtable}
\usepackage{booktabs}
\usepackage{array}
\usepackage{calc}
\usepackage{caption}
\newcounter{none}

% ── Math support ─────────────────────────────────────────────────────────────
\usepackage{amsmath}
\usepackage{amssymb}

% ── Make \paragraph and \subparagraph block-level ────────────────────────────
\usepackage{titlesec}
\titleformat{\paragraph}
  {\normalfont\normalsize\bfseries}{\theparagraph}{1em}{}
\titlespacing*{\paragraph}
  {0pt}{3.25ex plus 1ex minus .2ex}{1.5ex plus .2ex}
\titleformat{\subparagraph}
  {\normalfont\normalsize\bfseries}{\thesubparagraph}{1em}{}
\titlespacing*{\subparagraph}
  {0pt}{3.25ex plus 1ex minus .2ex}{1.5ex plus .2ex}

% ── Pandoc-generated tightlist support ───────────────────────────────────────
\providecommand{\tightlist}{%
  \setlength{\itemsep}{0pt}\setlength{\parskip}{0pt}}

% ── Pandoc 3.x figure sizing ────────────────────────────────────────────────
\newcommand{\pandocbounded}[1]{%
  \sbox0{#1}%
  \ifdim\wd0>\linewidth
    \resizebox{\linewidth}{!}{#1}%
  \else
    #1%
  \fi
}

% ── Unicode character fixes ──────────────────────────────────────────────────
\usepackage{newunicodechar}
\newunicodechar{✓}{$\checkmark$}
\newunicodechar{≠}{$\neq$}
\newunicodechar{≡}{$\equiv$}
\newunicodechar{≈}{$\approx$}
\newunicodechar{δ}{$\delta$}
\newunicodechar{−}{$-$}
\let\mystRightarrow\rightarrow
\renewcommand{\rightarrow}{\ensuremath{\mystRightarrow}}
% ─────────────────────────────────────────────────────────────────────────────

\hypersetup{colorlinks=true, allcolors=blue}

\title{Technical Summary: Necessary Constraints in Multi-Agent
Dissipative Optimization}
\author{Keith Lostracco}
\date{}

\begin{document}
\maketitle

\subsection{Model}\label{model}

Consider \(N \geq 2\) agents sharing a finite resource endowment
\(\mathbf{R} = (R_1, \ldots, R_n) \in \mathbb{R}_{>0}^n\). Each agent
\(i\) selects a strategy vector
\(\mathbf{x}_i \in \mathbb{R}_{\geq 0}^n\) and maximizes a utility
function \(U_i(\mathbf{x}_i)\) subject to shared resource constraints
\(\sum_i x_{ij} \leq R_j\) for all \(j\). Each \(U_i\) is \(C^2\),
strictly concave, monotonically increasing at the origin, and coercive
(\(U_i \to -\infty\) as \(\|\mathbf{x}_i\| \to \infty\)). These are
regularity conditions for Lagrange multiplier theory; the physical
content is diminishing returns and bounded metabolism.

Each agent maintains a boundary against environmental entropy at cost
\(C_{\text{maintain},i} = \gamma_i B_i\) (the baseline metabolic cost of
persistence).

Agent \(B\)'s boundary constraints are inequality constraints
\(g_{Bk}(\mathbf{x}_A) \leq 0\) on every other agent \(A\)'s problem,
where each \(g_{Bk}\) is continuously differentiable and convex,
satisfying \(g_{Bk}(\mathbf{0}) < 0\) (the zero-action strategy violates
no agent's boundary constraints). Agent \(A\)'s constrained problem is:

\[\max_{\mathbf{x}_A \geq \mathbf{0}} \; U_A(\mathbf{x}_A) \quad \text{s.t.} \quad x_{Aj} + \hat{x}_{-Aj} \leq R_j \;\; \forall j, \qquad g_{Bk}(\mathbf{x}_A) \leq 0 \;\; \forall B \neq A,\, k\]

The KKT stationarity condition at the optimum reads, at every resource
\(j\) that agent \(A\) consumes in positive quantity (\(x_{Aj}^* > 0\)):

\[\frac{\partial U_A}{\partial x_{Aj}}\bigg|_{\mathbf{x}_A^*} = \lambda_j + \sum_{B,k} \mu_{Bk} \frac{\partial g_{Bk}}{\partial x_{Aj}}\bigg|_{\mathbf{x}_A^*}\]

with the corresponding ``\(\leq\)'' at any corner \(x_{Aj}^* = 0\),
where the nonnegativity multiplier enters.

\subsection{Results}\label{results}

\textbf{Theorem 1 (Shadow Price).} Let
\(\mathbf{x}_A^*(\boldsymbol{c})\) denote agent \(A\)'s optimal strategy
as a function of the constraint parameters
\(\boldsymbol{c} = \{c_{Bk}\}\), where the \(k\)-th boundary constraint
of agent \(B\) is parameterized as \(g_{Bk}(\mathbf{x}_A) \leq c_{Bk}\).
Then:

\[\frac{\partial U_A^*}{\partial c_{Bk}} = \mu_{Bk}^* \geq 0\]

Every active boundary constraint has a quantifiable energetic cost to
the constrained agent.

\textbf{Lemma 1 (Inevitability of Collision).} Given a finite resource
endowment \(\mathbf{R}\) shared by \(N \geq 2\) agents with overlapping
resource needs, collision is guaranteed whenever \(N\) is sufficiently
large. Specifically, for any resource \(j\) essential to all agents,
collision occurs whenever
\(N \geq N_j^* := \lfloor R_j / \min_i x_{ij}^{\circ} \rfloor + 1\).

\textbf{Theorem 2 (Impossibility of Infinite Freedom).} Let \(N \geq 2\)
agents share a finite resource endowment \(\mathbf{R}\). Suppose each
agent has at least one essential resource (a resource \(j\) for which
\(\partial U_i / \partial x_{ij} |_{\mathbf{x}_i = \mathbf{0}} > 0\)) in
common with at least one other agent, and that aggregate unconstrained
demand exceeds supply for at least one resource
(\(\sum_i x_{ij}^{\circ} > R_j\) for some \(j\)). Then no feasible
strategy profile can simultaneously be unconstrained-optimal for all
agents:

\[\nexists \; (\mathbf{x}_1, \ldots, \mathbf{x}_N) \in \mathcal{F}_R \; \text{ such that } \; \mathbf{x}_i = \mathbf{x}_i^{\circ} \;\; \forall i\]

where \(\mathcal{F}_R\) is the resource-feasible set
(resource-conservation and nonnegativity constraints only).

\textbf{Corollary 2.1 (Necessity of Active Constraints).} In any
multi-agent system with resource collision, at least one constraint
restricting some agent must be active at the variational equilibrium; in
particular, in the resource-only GNEP over \(\mathcal{F}_R\) with
utilities separable across resources (as in the quadratic model), the
shared scarcity multiplier satisfies \(\lambda_j^* > 0\) for every
collided resource \(j\) (with binding boundary constraints, aggregate
consumption can be held strictly below \(R_j\), in which case
\(\lambda_j^* = 0\) and the constraint activity resides in the boundary
multipliers). Moreover, at any feasible profile at least one agent's
allocation of the collided resource lies strictly below its
unconstrained optimum. Some agent is therefore always constrained.

\textbf{Theorem 3 (Variational Equilibrium).} Under the regularity
conditions, the GNEP possesses a Variational Equilibrium
\((\mathbf{x}_1^*, \ldots, \mathbf{x}_N^*, \boldsymbol{\lambda}^*)\) in
which:

\begin{enumerate}
\def\labelenumi{(\alph{enumi})}
\item
  The shared shadow prices \(\lambda_j^*\) are identical for all agents.
\item
  The stationarity condition for each agent \(i\) reads, at every
  resource \(j\) with \(x_{ij}^* > 0\):
\end{enumerate}

\[\frac{\partial U_i}{\partial x_{ij}}\bigg|_{\mathbf{x}_i^*} = \lambda_j^* + \sum_{B \neq i, k} \mu_{Bk}^{(i)} \frac{\partial g_{Bk}}{\partial x_{ij}}\bigg|_{\mathbf{x}_i^*}\]

(with the corresponding ``\(\leq\)'' at any corner \(x_{ij}^* = 0\),
where the nonnegativity multiplier enters).

\begin{enumerate}
\def\labelenumi{(\alph{enumi})}
\setcounter{enumi}{2}
\tightlist
\item
  No agent can improve its utility by unilaterally deviating from
  \(\mathbf{x}_i^*\), given the constraints.
\end{enumerate}

\textbf{Corollary 3.1 (Symmetry of Scarcity).} At the variational
equilibrium, all agents experience the same marginal cost
\(\lambda_j^*\) per unit of each shared resource.

\textbf{Theorem 4 (Welfare Optimality).} Under the regularity
conditions, the variational equilibrium is the unique maximizer of total
social welfare \(W = \sum_i U_i(\mathbf{x}_i)\) over the feasible set
\(\mathcal{F}\) (shared resource and boundary constraints):

\[\sum_{i=1}^{N} U_i(\mathbf{x}_i^*) > \sum_{i=1}^{N} U_i(\mathbf{x}_i) \qquad \forall \; (\mathbf{x}_1, \ldots, \mathbf{x}_N) \in \mathcal{F} \setminus \{\mathbf{x}^*\}\]

\textbf{Corollary 4.1 (Welfare Dominance).} The variational equilibrium
is welfare-superior to every other allocation in \(\mathcal{F}\); any
feasible greedy or equal-split allocation differing from the VE is
strictly dominated. (Greedy dominance requires the greedy allocation to
be feasible; for symmetric agents the VE coincides with the equal
split.)

\textbf{Corollary 4.2 (Constraint Value Scales with Scarcity).} In the
symmetric quadratic model, on the window
\(\alpha/\beta \leq R < 2\alpha/\beta\) (the portion of the collision
regime in which the defector's unconstrained optimum remains feasible),
the VE-vs-defection welfare gap
\(V(R) = \tfrac{\beta}{4}(2\alpha/\beta - R)^2\) is strictly increasing
in scarcity; under abundance the constraint system is slack (\(V = 0\)),
and under extreme scarcity (\(R < \alpha/\beta\)) the defector absorbs
the entire endowment and \(V(R) = \tfrac{\beta}{4}R^2\), increasing in
\(R\). The general case is left as a conjecture.

\textbf{Proposition 1 (Unconstrained Regime Implies Conflict).} If all
boundary constraints are removed and the system is in collision, then:
(a) no stable allocation exists without an enforcement mechanism; (b)
each agent's attempt to realize its unconstrained optimum generates a
negative externality on other agents; (c) the resulting contested
resource must be resolved by direct energy expenditure.

\subsection{Notation}\label{notation}

{\def\LTcaptype{none} % do not increment counter
\begin{longtable}[]{@{}
  >{\raggedright\arraybackslash}p{(\linewidth - 2\tabcolsep) * \real{0.5000}}
  >{\raggedright\arraybackslash}p{(\linewidth - 2\tabcolsep) * \real{0.5000}}@{}}
\toprule\noalign{}
\begin{minipage}[b]{\linewidth}\raggedright
Symbol
\end{minipage} & \begin{minipage}[b]{\linewidth}\raggedright
Meaning
\end{minipage} \\
\midrule\noalign{}
\endhead
\bottomrule\noalign{}
\endlastfoot
\(N\) & Number of agents \\
\(n\) & Number of resource dimensions \\
\(\mathbf{x}_i \in \mathbb{R}_{\geq 0}^n\) & Agent \(i\)'s strategy
(resource-allocation) vector \\
\(R_j\) & Total available quantity of resource \(j\) \\
\(U_i(\mathbf{x}_i)\) & Agent \(i\)'s utility (boundary-maintenance
objective) \\
\(\alpha_{ij}, \beta_{ij}\) & Marginal yield and diminishing-returns
parameters (quadratic model) \\
\(\mathbf{x}_i^{\circ}\) & Agent \(i\)'s unconstrained optimum \\
\(g_{Bk}(\mathbf{x}_A)\) & Constraint function encoding agent \(B\)'s
\(k\)-th boundary constraint \\
\(\lambda_j\) & Lagrange multiplier for resource \(j\) scarcity \\
\(\mu_{Bk}\) & Lagrange multiplier (shadow price) for agent \(B\)'s
\(k\)-th boundary constraint \\
\(\mathcal{F}_R\) & Resource-feasible set: resource-conservation and
nonnegativity constraints only (Theorem 2) \\
\(\mathcal{F}\) & Feasible strategy set under resource \emph{and}
boundary constraints (Theorem 4) \\
\end{longtable}
}

\end{document}
